\documentclass[12pt,a4paper]{article}
\usepackage[utf8]{inputenc}
\usepackage[T1]{fontenc}
\usepackage{geometry}
\usepackage{longtable}
\usepackage{hyperref}
\geometry{a4paper, margin=2.5cm}
\title{\textbf{Manual do Usuário}\\Sistema Reserva de Salas de Estudo}
\author{Ailton Baren Pereira da Silva}
\date{Junho de 2026}
\begin{document}
\maketitle

\section{Acesso}
A aplicação deve ser acessada pelo navegador no endereço:
\begin{verbatim}
http://127.0.0.1:5000
\end{verbatim}

\section{Usuários Iniciais}
\begin{longtable}{|p{4cm}|p{6cm}|p{4cm}|}
\hline
\textbf{Perfil} & \textbf{E-mail} & \textbf{Senha} \\ \hline
Administrador & admin@uema.com & admin123 \\ \hline
Aluno & aluno@uema.com & aluno123 \\ \hline
Aluno & aluno1@uema.com & aluno123 \\ \hline
Aluno & aluno2@uema.com & aluno123 \\ \hline
Aluno & aluno3@uema.com & aluno123 \\ \hline
\end{longtable}

\section{Aluno}
\subsection{Dashboard}
Após o login, o aluno visualiza a quantidade de reservas ativas, próximas reservas, salas disponíveis e botão para consultar horários.

\subsection{Consultar Salas}
No menu superior, acesse \textbf{Salas} para visualizar nome, localização, capacidade e status das salas. Quando uma sala possuir bloqueio ativo para o dia atual, ela será exibida como \textbf{Bloqueada hoje} e não aparecerá no filtro de salas disponíveis.

\subsection{Consultar Horários}
No menu superior, acesse \textbf{Horários}. Selecione sala e data. O sistema apresenta horários disponíveis, ocupados e bloqueados em blocos de uma hora.

\subsection{Criar Reserva}
Na tela de horários, informe a finalidade da reserva e clique no botão do horário desejado. O sistema valida conflitos, data passada, duração fixa de uma hora, sala inativa, sala bloqueada e limite de reservas ativas.

\subsection{Minhas Reservas}
No menu \textbf{Minhas reservas}, o aluno consulta reservas ativas e canceladas.

\subsection{Cancelar Reserva}
Para cancelar, clique em \textbf{Cancelar}, informe o motivo e confirme a ação. O motivo fica registrado.

\section{Administrador}
\subsection{Painel Administrativo}
O painel exibe indicadores de usuários, salas, salas bloqueadas e reservas ativas.

\subsection{Gerenciar Usuários}
O administrador pode cadastrar novos alunos e administradores, além de editar usuários existentes, informando nome, e-mail, matrícula, perfil, senha e status ativo ou inativo. Na edição, senha em branco mantém a senha atual.

\subsection{Gerenciar Salas}
O administrador pode cadastrar, editar, ativar, desativar, bloquear e excluir salas. A exclusão é permitida apenas para salas sem reservas ou bloqueios registrados; caso exista histórico, a orientação é desativar a sala para impedir novas reservas sem apagar registros antigos.

\subsection{Reservas Administrativas}
O administrador pode visualizar todas as reservas, filtrar por aluno, sala, status e data, além de cancelar reservas com motivo.

\subsection{Bloqueios}
O administrador pode bloquear salas por período, informando sala, data inicial, data final e motivo. Também pode cancelar bloqueios.

\subsection{Configurações}
O administrador pode configurar o limite máximo de reservas ativas por aluno.

\section{Mensagens}
O sistema exibe mensagens de sucesso e erro em português para login, cadastros, reservas, bloqueios, cancelamentos e validações.

\end{document}
